% Einheit 1 von 2 – Streifen- und Kreisdiagramm (Katalog: daten.md, Einheit 3)
\setcounter{aufgabe}{11}
\zweigkopf{e1}{Einheit 1 von 2 · Streifen- und Kreisdiagramm}{Hier lernst du, Anteile als Streifen und als Kreisdiagramm darzustellen · neu oder schon bekannt – je nach Buch (am Gymnasium schon seit Klasse 6) · P10 oft · baut auf: Häufigkeiten und Anteile (Kennst du schon), Prozentsatz, Längen in cm und mm, Winkel zeichnen}

\begin{aufgabe}{\textbf{Ich erkenne, was das Ganze ist.} Unterstreiche in jedem Satz die Zahl, die für das Ganze steht. Rechne nichts aus.}
\begin{teile}
\teil Von den 28 Kindern der Klasse 7b kommen 12 mit dem Bus.
\teil Bei einer Umfrage unter 250 Jugendlichen nannten 85 Fußball als Lieblingssport.
\teil 45 der 180 Bäume im Stadtpark sind Eichen.
\teil Ein Sportverein hat 640 Mitglieder, 160 davon sind jünger als 18 Jahre.
\teil Bei der Klassensprecherwahl stimmten 24 Kinder ab: Mia bekam 11 Stimmen, Tom 9 Stimmen, und 4 Zettel waren ungültig.
\end{teile}
\end{aufgabe}

\begin{aufgabe}{\textbf{Ich erkenne, ob eine Zahl ein Anteil in Prozent oder ein Winkel in Grad ist.} Die Sätze gehören zu einem Kreisdiagramm „Wie kommst du zur Schule?“; die Einheit fehlt. Kreuze an, was die Zahl ist.}
\begin{teile}
\teil Der ganze Kreis hat 360. \hfill \kreuz{Anteil in \%}\quad\kreuz{Winkel in °}
\teil Ein Viertel aller Kinder fährt Rad, das sind 25. \hfill \kreuz{Anteil in \%}\quad\kreuz{Winkel in °}
\teil Lena misst den Sektor „Bus“ an der Kreismitte: 144. \hfill \kreuz{Anteil in \%}\quad\kreuz{Winkel in °}
\teil Alle Anteile zusammen ergeben 100. \hfill \kreuz{Anteil in \%}\quad\kreuz{Winkel in °}
\teil Die Hälfte der Kinder kommt zu Fuß, also 50. \hfill \kreuz{Anteil in \%}\quad\kreuz{Winkel in °}
\teil Der Sektor „Auto“ ist ein rechter Winkel, also 90. \hfill \kreuz{Anteil in \%}\quad\kreuz{Winkel in °}
\end{teile}
\end{aufgabe}

\begin{aufgabe}{\textbf{Ich kann einen Anteil als Streifen darstellen.} Der Streifen ist 10\,cm lang und steht für $100\,\%$. Gib an, wie lang der Abschnitt ist, und färbe ihn im Streifen.}
Beispiel: $30\,\%$ von 10\,cm sind 3\,cm.

\streifen{30}{$0\,\%$}{$100\,\%$}

\begin{teile}
\teil $40\,\%$\par\noindent\begin{minipage}[b]{11cm}\streifenleer\end{minipage}\hfill\leerfeld[cm]
\teil $70\,\%$\par\noindent\begin{minipage}[b]{11cm}\streifenleer\end{minipage}\hfill\leerfeld[cm]
\teil $20\,\%$\par\noindent\begin{minipage}[b]{11cm}\streifenleer\end{minipage}\hfill\leerfeld[cm]
\teil $90\,\%$\par\noindent\begin{minipage}[b]{11cm}\streifenleer\end{minipage}\hfill\leerfeld[cm]
\end{teile}
\end{aufgabe}

\begin{aufgabe}{\textbf{Ich kann einen Anteil als Streifen darstellen – weiter.}}
\begin{teile}
\teil $65\,\%$\par\noindent\begin{minipage}[b]{11cm}\streifenleer\end{minipage}\hfill\leerfeld[cm]
\teil Zeichne einen Streifen mit drei Abschnitten und beschrifte sie: Rad $45\,\%$, Bus $40\,\%$, zu Fuß $15\,\%$.\par
\streifenleer
\teil Beim Frühstück: Müsli $50\,\%$, Brot $20\,\%$, der Rest frühstückt nicht. Ergänze den fehlenden Anteil und zeichne den Streifen mit allen drei Abschnitten.\\ kein Frühstück: \leerfeld[\%]\par
\streifenleer
\teil Beim Pausenverkauf der Schule gingen weg: Brötchen $40\,\%$, Getränke $30\,\%$, Müsliriegel $25\,\%$, Kekse $5\,\%$. Stelle die Anteile in einem Streifendiagramm dar und beschrifte die Abschnitte. (P10 2018)\par
\streifenleer
\end{teile}
\end{aufgabe}

\begin{aufgabe}{\textbf{Ich kann ausrechnen, welcher Winkel im Kreisdiagramm zu einem Anteil gehört.} Der ganze Kreis ($360^\circ$) steht für $100\,\%$. Berechne den Winkel.}
Beispiel: $20\,\% \to 0{,}2 \cdot 360^\circ = 72^\circ$
\begin{teilezwei}
\tz $50\,\%$: \leerfeld[$^\circ$] & \tz $25\,\%$: \leerfeld[$^\circ$] \\
\tz $10\,\%$: \leerfeld[$^\circ$] & \tz $75\,\%$: \leerfeld[$^\circ$] \\
\tz $15\,\%$: \leerfeld[$^\circ$] & \tz $5\,\%$: \leerfeld[$^\circ$] \\
\tz $12{,}5\,\%$: \leerfeld[$^\circ$] & \\
\end{teilezwei}
Berechne erst den Anteil, dann den Winkel.
\begin{teile}
\teil $\frac{3}{8}$ der Kinder spielen ein Instrument. \hfill \leerfeld[$^\circ$]
\teil 13 von 40 Kindern haben ein Haustier. \hfill \leerfeld[$^\circ$]
\teil Von den 13 Mitgliedern einer Theater-AG sind 4 Jungen. Berechne den Winkel für die Jungen und runde auf ganze Grad. (P10 2025) \hfill \leerfeld[$^\circ$]
\end{teile}
\end{aufgabe}

\begin{aufgabe}{\textbf{Ich kann ein Kreisdiagramm zeichnen (kennst du wahrscheinlich seit Klasse 6).} Berechne die Winkel, trage sie im Kreis ab dem Radius nach oben mit dem Geodreieck an und beschrifte jeden Sektor.}
Beispiel: Rad $25\,\% \to 0{,}25 \cdot 360^\circ = 90^\circ$; Radius nach oben, $90^\circ$ an der Kreismitte antragen, Sektor mit „Rad $25\,\%$“ beschriften.

\noindent\begin{minipage}[t]{0.48\linewidth}
\begin{teile}
\teil Ein Sektor für $40\,\%$.
\end{teile}
\kreisleer[2.3]
\end{minipage}\hfill
\begin{minipage}[t]{0.48\linewidth}
\begin{teile}
\teil Tee $50\,\%$, Kakao $45\,\%$, Milch $5\,\%$.
\end{teile}
\kreisleer[2.3]
\end{minipage}
\end{aufgabe}

\begin{aufgabe}{\textbf{Ich kann ein Kreisdiagramm zeichnen – weiter.} In einer Umfrage zur Lieblingsbeschäftigung in der Freizeit nannten $42{,}5\,\%$ Sport, $27{,}5\,\%$ Musik, $18{,}5\,\%$ Lesen, der Rest Sonstiges.}
\noindent\begin{minipage}[t]{0.5\linewidth}
\begin{teile}
\teil Ergänze die Tabelle. Runde die Winkel auf ganze Grad.
\end{teile}
\sachtabelle{lcc}{Freizeit & Anteil & Winkel}{Sport & $42{,}5\,\%$ & \leerzelle \\ Musik & $27{,}5\,\%$ & \leerzelle \\ Lesen & $18{,}5\,\%$ & \leerzelle \\ Sonstiges & \leerzelle & \leerzelle}
\begin{teile}
\teil Zeichne das Kreisdiagramm und beschrifte die Sektoren. (P10 2017)
\end{teile}
\end{minipage}\hfill
\begin{minipage}[t]{0.45\linewidth}
\vspace{0pt}\kreisleer[2.5]
\end{minipage}
\end{aufgabe}

\begin{aufgabe}{\textbf{Ich kann schätzen, welcher Anteil zu einem Sektor gehört.} Schätze für jeden Sektor, ob er die Hälfte, ein Drittel, ein Viertel oder ein Sechstel des Kreises ist, und gib den Anteil ungefähr in Prozent an.}
\noindent\begin{minipage}[t]{0.48\linewidth}
\centering\kreisdiagramm[1.6]{A/50, B/25, C/25}
\end{minipage}\hfill
\begin{minipage}[t]{0.48\linewidth}
\centering\kreisdiagramm[1.6]{D/33.333, E/50, F/16.667}
\end{minipage}
\begin{teile}
\teil A: \leerfeld \quad ungefähr \leerfeld[\%]
\teil B: \leerfeld \quad ungefähr \leerfeld[\%]
\teil D: \leerfeld \quad ungefähr \leerfeld[\%]
\teil F: \leerfeld \quad ungefähr \leerfeld[\%]
\end{teile}
\end{aufgabe}

\begin{aufgabe}{\textbf{Ich kann die Sektoren eines Kreisdiagramms den Anteilen zuordnen.} Die Sektoren sind nicht beschriftet. Schreibe an jeden Strich, welcher Anteil dazugehört.}
\noindent\begin{minipage}[t]{0.40\linewidth}
\begin{teile}
\teil Bücher $55\,\%$, Handy $25\,\%$, Fernsehen $20\,\%$
\end{teile}
\kreisdiagrammleer{a/25, b/55, c/20}
\end{minipage}\hfill
\begin{minipage}[t]{0.56\linewidth}
\begin{teile}
\teil Familie Kaya hat in einem Jahr 3200\,kWh Strom verbraucht: Unterhaltung $34\,\%$, Kühlen $24\,\%$, Kochen $18\,\%$, Waschen $10\,\%$, der Rest für Sonstiges.
\end{teile}
\kreisdiagrammleer{a/34, b/10, c/24, d/14, e/18}
\end{minipage}
\begin{teile}
\teil Wie viel Prozent entfallen bei Familie Kaya auf Sonstiges? \leerfeld[\%]
\teil Zum Laden der Handys hat die Familie 72\,kWh gebraucht. Berechne den Winkel des Sektors, der dazu gehören würde. Runde auf eine Nachkommastelle. (P10 2024)\\ \leerfeld[$^\circ$]
\end{teile}
\end{aufgabe}

\begin{aufgabe}{\textbf{Ich finde den Fehler beim Winkel im Kreisdiagramm.} Jana hat für eine Umfrage (Rad $40\,\%$, Bus $45\,\%$, zu Fuß $15\,\%$) die Winkel so berechnet:}
\rechnung{\text{Rad: } 40\,\% \text{ von } 360^\circ &= 0{,}4 \cdot 360^\circ = 144^\circ \\ \text{Bus: } 45\,\% \text{ von } 360^\circ &= 0{,}45 \cdot 360^\circ = 162^\circ \\ \text{zu Fuß: } 15\,\% \text{ von } 360^\circ &= 15^\circ}
\begin{teile}
\teil Finde den Fehler, benenne ihn und korrigiere die Rechnung.
\teil Berechne die Winkel für Kakao $55\,\%$, Tee $20\,\%$ und Wasser $25\,\%$.
\end{teile}
Timo soll den Winkel für 12 von 300 befragten Kindern berechnen. Er rechnet so:
\rechnung{12 \text{ von } 300 &\to 12\,\% \text{ von } 360^\circ \\ &= 0{,}12 \cdot 360^\circ = 43{,}2^\circ}
\begin{teile}
\teil Finde den Fehler, benenne ihn und korrigiere die Rechnung.
\teil Berechne den Winkel für 9 von 450 befragten Kindern. \leerfeld[$^\circ$]
\end{teile}
\end{aufgabe}

\begin{aufgabe}{\textbf{Ich kann begründen, wie Anteil und Winkel im Kreisdiagramm zusammenhängen.} Antworte ohne Taschenrechner.}
\begin{teile}
\teil Begründe, warum $1\,\%$ im Kreisdiagramm genau $3{,}6^\circ$ entspricht.
\teil In der 7a (20 Kinder) und in der 7b (30 Kinder) fahren jeweils $50\,\%$ der Kinder mit dem Rad. Anna meint: „Im Kreisdiagramm der 7b muss der Sektor für Rad größer sein, weil dort mehr Kinder Rad fahren.“ Hat Anna recht? Begründe.
\end{teile}
\end{aufgabe}

\begin{aufgabe}{\textbf{Ich kann entscheiden, ob ein Kreisdiagramm zu einer Umfrage passt.} Die Schülerzeitung hat 80 Kinder gefragt: „Welchen Sport machst du in der Freizeit?“ Jedes Kind durfte mehrere Sportarten nennen.}

\sachtabelle{lcc}{Sport & Nennungen & Anteil}{Fußball & 36 & \leerzelle \\ Schwimmen & 24 & \leerzelle \\ Radfahren & 32 & \leerzelle \\ Tanzen & 16 & \leerzelle}

\begin{teile}
\teil Berechne für jede Sportart den Anteil an allen 80 Kindern in Prozent und trage ihn ein.
\teil Addiere die Anteile. \leerfeld[\%]
\teil Die Redaktion will die Ergebnisse als Kreisdiagramm zeigen. Geht das? Begründe und schlage eine passende Darstellung vor.
\end{teile}
\end{aufgabe}
